\documentclass[handout]{beamer}
\mode<presentation>
\usepackage{xcolor}
\usepackage[toc]{multitoc}
\usepackage{color}
\usepackage{graphicx}
\graphicspath{{../buku_ajar/buku_ajar_logika_matematika/figures/}}
\usepackage{tikz}
\usetikzlibrary{shapes,arrows,positioning,calc}
\usepackage{listings}
\usepackage{lmodern}
\usetheme[secheader]{Boadilla}
\usecolortheme{whale}
\setbeamertemplate{caption}[numbered]
\setbeamertemplate{bibliography item}[text]
\setbeamertemplate{navigation symbols}{}

\theoremstyle{definition}
\newtheorem{definisi}[theorem]{Definisi}
\theoremstyle{example}
\newtheorem{contoh}[theorem]{Contoh}

\renewcommand*{\multicolumntoc}{2}

\setbeamertemplate{footline}
{
  \leavevmode%
  \hbox{%
  \begin{beamercolorbox}[wd=.333333\paperwidth,ht=2.25ex,dp=1ex,center]{author in head/foot}%
    \usebeamerfont{author in head/foot}\insertshortauthor%
  \end{beamercolorbox}%
  \begin{beamercolorbox}[wd=.433333\paperwidth,ht=2.25ex,dp=1ex,center]{title in head/foot}%
    \usebeamerfont{title in head/foot}\insertshorttitle
  \end{beamercolorbox}%
  \begin{beamercolorbox}[wd=.233333\paperwidth,ht=2.25ex,dp=1ex,right]{date in head/foot}%
    \usebeamerfont{date in head/foot}\insertshortdate{}\hspace*{2em} \insertframenumber{} / \inserttotalframenumber\hspace*{2ex} 
  \end{beamercolorbox}}%
  \vskip0pt%
}

\subtitle[Logika Matematika]{Logika Matematika}
\title[Aturan Inferensi]{Aturan Inferensi dan Penarikan Kesimpulan}
\author{Cecep Suwanda, S.Si., M.Kom.}
\date{Maret 2025}
\institute{Universitas Bale Bandung}

\begin{document}

\maketitle

\section*{Daftar Isi}
\begin{frame}{Daftar Isi}
\frametitle{Daftar Isi}
\tableofcontents
\end{frame}

% ============================================================
% SECTION 5-1: Pengantar Inferensi
% ============================================================
\section{Pengantar Inferensi}

\begin{frame}
\frametitle{Inferensi}
\begin{definisi}[Inferensi Logika]
Mekanisme deduktif formal yang memanfaatkan aturan matematika untuk memperoleh konklusi baru dari himpunan premis.
\end{definisi}
\begin{block}{Jaminan}
Apabila semua premis benar, maka konklusi yang dihasilkan juga benar.
\end{block}
\end{frame}

\begin{frame}
\frametitle{Struktur Argumen Formal}
\begin{center}
\begin{tabular}{ll}
$P_1$ & (Premis pertama) \\
$P_2$ & (Premis kedua) \\
$\vdots$ & \\
$P_n$ & (Premis ke-$n$) \\
\hline
$\therefore Q$ & (Konklusi)
\end{tabular}
\end{center}
\vspace{0.3cm}
Secara horizontal: $(P_1 \wedge P_2 \wedge \ldots \wedge P_n) \rightarrow Q$
\end{frame}

\begin{frame}
\frametitle{Validitas Argumen}
\begin{block}{Teorema}
Argumen dengan premis $P_1, \ldots, P_n$ dan konklusi $Q$ disebut \textbf{valid} jika dan hanya jika
\[ [P_1 \wedge P_2 \wedge \ldots \wedge P_n] \rightarrow Q \]
merupakan \textbf{tautologi}.
\end{block}
\begin{block}{}
Jika ada interpretasi: semua premis benar tetapi $Q$ salah $\Rightarrow$ argumen \textbf{tidak valid} (fallacy).
\end{block}
\end{frame}

% ============================================================
% SECTION 5-2: Modus Ponens dan Modus Tollens
% ============================================================
\section{Modus Ponens dan Modus Tollens}

\begin{frame}
\frametitle{Modus Ponens}
\begin{block}{Skema}
\begin{tabular}{cl}
$p \rightarrow q$ & (Implikasi) \\
$p$ & (Anteseden benar) \\
\hline
$\therefore q$ & (Konsekuen benar)
\end{tabular}
\end{block}
\begin{contoh}
$p \rightarrow q$: ``Jika Rini skor TOEFL 500, maka memenuhi yudisium.''\\
$p$: ``Rini skor TOEFL 500.''\\
$\therefore q$: ``Rini memenuhi yudisium.''
\end{contoh}
\end{frame}

\begin{frame}
\frametitle{Modus Tollens}
\begin{block}{Skema}
\begin{tabular}{cl}
$p \rightarrow q$ & (Implikasi) \\
$\neg q$ & (Konsekuen salah) \\
\hline
$\therefore \neg p$ & (Anteseden salah)
\end{tabular}
\end{block}
\begin{contoh}
$p \rightarrow q$: ``Jika penjaga kolam pelaku, maka sepatunya basah lumpur.''\\
$\neg q$: ``Sepatu penjaga kering.''\\
$\therefore \neg p$: ``Penjaga kolam bukan pelaku.''
\end{contoh}
\end{frame}

% ============================================================
% SECTION 5-3: Silogisme
% ============================================================
\section{Silogisme Hipotetis dan Disjungtif}

\begin{frame}
\frametitle{Silogisme Hipotetis}
\begin{block}{Skema (sifat transitif implikasi)}
\begin{tabular}{c}
$p \rightarrow q$ \\
$q \rightarrow r$ \\
\hline
$\therefore p \rightarrow r$
\end{tabular}
\end{block}
\begin{contoh}
``Jika belajar logika $\Rightarrow$ paham aljabar.''\\
``Jika paham aljabar $\Rightarrow$ lulus ujian.''\\
$\therefore$ ``Jika belajar logika $\Rightarrow$ lulus ujian.''
\end{contoh}
\end{frame}

\begin{frame}
\frametitle{Silogisme Disjungtif}
\begin{block}{Skema}
\begin{tabular}{ccc}
$p \vee q$ & \quad & $p \vee q$ \\
$\neg p$ & atau & $\neg q$ \\
\hline
$\therefore q$ & & $\therefore p$
\end{tabular}
\end{block}
\begin{contoh}
$p \vee q$: ``Pelaku lari lewat jendela ATAU plafon.''\\
$\neg p$: ``Tidak ada bukti lewat jendela.''\\
$\therefore q$: ``Pelaku lari lewat plafon.''
\end{contoh}
\end{frame}

% ============================================================
% SECTION 5-4: Konjungsi, Penambahan, Simplifikasi
% ============================================================
\section{Aturan Konjungsi, Penambahan, Simplifikasi}

\begin{frame}
\frametitle{Aturan Konjungsi}
\begin{block}{Skema}
\begin{tabular}{c}
$p$ \\
$q$ \\
\hline
$\therefore p \wedge q$
\end{tabular}
\end{block}
\begin{contoh}
$p$: Jakarta ibu kota Indonesia.\\
$q$: Bangkok ibu kota Thailand.\\
$\therefore p \wedge q$: Jakarta ibu kota Indonesia DAN Bangkok ibu kota Thailand.
\end{contoh}
\end{frame}

\begin{frame}
\frametitle{Aturan Penambahan (Addition)}
\begin{block}{Skema}
\begin{tabular}{c}
$p$ \\
\hline
$\therefore p \vee q$
\end{tabular}
\end{block}
\begin{contoh}
$p$: ``Budi berolahraga setiap pagi.''\\
$\therefore p \vee q$: ``Budi berolahraga setiap pagi ATAU hujan di kutub selatan.''
\end{contoh}
\end{frame}

\begin{frame}
\frametitle{Aturan Simplifikasi}
\begin{block}{Skema}
\begin{tabular}{ccc}
$p \wedge q$ & \quad atau \quad & $p \wedge q$ \\
\hline
$\therefore p$ & & $\therefore q$
\end{tabular}
\end{block}
\begin{contoh}
$p \wedge q$: ``Air mendidih 100°C DAN api membakar.''\\
$\therefore p$: ``Air mendidih 100°C.'' \quad atau \quad $\therefore q$: ``Api membakar.''
\end{contoh}
\end{frame}

% ============================================================
% SECTION 5-5: Aturan Resolusi
% ============================================================
\section{Aturan Resolusi}

\begin{frame}
\frametitle{Prinsip Resolusi}
\begin{block}{Mekanisme}
Dua klausa disjungsi dengan literal komplementer ($p$ dan $\neg p$). Eliminasi pasangan $\Rightarrow$ resolvent.
\end{block}
\begin{block}{Skema}
\begin{tabular}{cl}
$p \vee q$ & (Klausa 1) \\
$\neg p \vee r$ & (Klausa 2) \\
\hline
$\therefore q \vee r$ & (Resolvent)
\end{tabular}
\end{block}
\end{frame}

\begin{frame}
\frametitle{Contoh Resolusi}
\begin{contoh}
$p \vee q$: ``Andi di Surabaya ATAU bertugas di Jakarta.''\\
$\neg p \vee r$: ``Andi tidak di Surabaya ATAU mobilnya di jalan tol.''\\
$\therefore q \vee r$: ``Andi bertugas di Jakarta ATAU mobilnya di jalan tol.''
\end{contoh}
\begin{block}{}
Dasar inferensi pada PROLOG dan sistem berbasis logika.
\end{block}
\end{frame}

% ============================================================
% SECTION 5-6: Fallacy
% ============================================================
\section{Kesesatan Berpikir (Fallacy)}

\begin{frame}
\frametitle{Fallacy of Affirming the Consequent}
\begin{block}{Skema yang TIDAK VALID}
\begin{tabular}{cl}
$p \rightarrow q$ & (Premis) \\
$q$ & (Premis) \\
\hline
$\therefore p$ & (Kesimpulan salah)
\end{tabular}
\end{block}
\begin{contoh}
$p \rightarrow q$: ``Jika kompor meledak $\Rightarrow$ gedung terbakar.''\\
$q$: ``Gedung terbakar.''\\
$\therefore p$ (salah): ``Pasti kompor meledak.''\\
\textit{Kebakaran bisa dari faktor lain.}
\end{contoh}
\end{frame}

\begin{frame}
\frametitle{Fallacy of Denying the Antecedent}
\begin{block}{Skema yang TIDAK VALID}
\begin{tabular}{cl}
$p \rightarrow q$ & (Premis) \\
$\neg p$ & (Premis) \\
\hline
$\therefore \neg q$ & (Kesimpulan salah)
\end{tabular}
\end{block}
\begin{contoh}
$p \rightarrow q$: ``Jika dapat bonus $\Rightarrow$ tabungan bertambah.''\\
$\neg p$: ``Tidak dapat bonus.''\\
$\therefore \neg q$ (salah): ``Tabungan tidak bertambah.''\\
\textit{Tabungan bisa bertambah dari sumber lain.}
\end{contoh}
\end{frame}

% ============================================================
% SECTION 5-7: Bukti Formal Berantai
% ============================================================
\section{Bukti Formal Berantai}

\begin{frame}
\frametitle{Prosedur Bukti Formal}
\begin{itemize}
  \item Setiap langkah harus memiliki justifikasi: nomor baris + nama aturan inferensi.
  \item Metode: \textbf{deduksi formal} (formal proof of validity).
\end{itemize}
\end{frame}

\begin{frame}
\frametitle{Contoh: Bukti Berbaris}
Premis: (1) $\neg p \rightarrow q$, (2) $\neg q$, (3) $\neg q \rightarrow r$. Konklusi: $r \wedge p$
\begin{center}
\small
\begin{tabular}{|c|l|l|}
\hline
\textbf{Baris} & \textbf{Ekspresi} & \textbf{Justifikasi} \\
\hline
1. & $\neg p \rightarrow q$ & Premis 1 \\
2. & $\neg q$ & Premis 2 \\
3. & $\neg q \rightarrow r$ & Premis 3 \\
\hline
4. & $\neg(\neg p)$ & Modus Tollens (1,2) \\
5. & $p$ & Involusi (4) \\
6. & $r$ & Modus Ponens (3,2) \\
\hline
7. & $\therefore r \wedge p$ & Konjungsi (5,6) \\
\hline
\end{tabular}
\end{center}
\end{frame}

% ============================================================
% RANGKUMAN
% ============================================================
\section{Rangkuman}

\begin{frame}
\frametitle{Rangkuman}
\begin{itemize}
  \item Argumen valid $\Leftrightarrow$ $(P_1 \wedge \ldots \wedge P_n) \rightarrow Q$ tautologi.
  \item Modus Ponens: $p \rightarrow q$, $p$ $\vdash$ $q$; Modus Tollens: $p \rightarrow q$, $\neg q$ $\vdash$ $\neg p$.
  \item Silogisme Hipotetis: $p \rightarrow q$, $q \rightarrow r$ $\vdash$ $p \rightarrow r$.
  \item Silogisme Disjungtif: $p \vee q$, $\neg p$ $\vdash$ $q$.
  \item Konjungsi, Penambahan, Simplifikasi, Resolusi.
  \item Fallacy: Affirming Consequent, Denying Antecedent.
  \item Bukti formal berantai mensyaratkan justifikasi tiap langkah.
\end{itemize}
\end{frame}

\begin{frame}
\frametitle{Kata Kunci}
\begin{center}
Inferensi $\bullet$ Validitas $\bullet$ Modus Ponens $\bullet$ Modus Tollens $\bullet$ Silogisme\\
Resolusi $\bullet$ Fallacy $\bullet$ Bukti Formal Berantai $\bullet$ Deduksi
\end{center}
\end{frame}

\end{document}
